\section{Appendix 20: Scenario Upsilon - Shock}

\subsection{The Legacy Paradigm \\\& Its Failures}
In legacy systems, sudden infrastructure or economic shocks result in immediate denial of service. Cloud-dependent applications freeze, centralized databases become unreachable, and digital interfaces turn into infinite loading spinners. This brittle architecture leads to information vacuums, widespread panic, and the complete failure of logistics when they are needed most.

\subsection{The Incremental Automation Pathway}
Transitioning to resilient shock-handling requires three phases:
Phase 1: Manual human entry of local resource states into offline-first digital forms.
Phase 2: Ambient sensor networks automatically updating resource manifests via short-range protocols (Bluetooth/LoRa), providing AR overlays of available supplies.
Phase 3: Fully ephemeral interactions where the system auto-reallocates critical resources based on pre-computed CRDTs, requiring zero human data entry.

\subsection{The Human Boundary (Limits of Automation)}
While the system can automate the cataloging and allocation of baseline resources, the interface must enforce a hard stop for life-threatening triage decisions. Human operators must manually confirm the diversion of critical power or medical supplies, ensuring that algorithmic triage does not silently override human ethics. Friction is intentionally introduced to demand explicit consent.

\subsection{Interface Mechanics (The "How")}
The interface relies on offline CRDT synchronization to ensure every node has a reasonably consistent view of the world even in a partitioned network. As thermodynamic limits are reached, the UI dynamically degrades from rich AR overlays to high-contrast, low-power E-ink displays.
